% =============================================
% File: chapters/experiment_results.tex
% Chương 4: Kết quả thực nghiệm
% =============================================

\chapter{KẾT QUẢ THỰC NGHIỆM}
\label{chap:ketqua}

Chương này trình bày kết quả cụ thể của từng giai đoạn thực nghiệm đã được mô tả
trong Chương~\ref{chap:phuongphap}: hàm truyền động cơ nhận dạng được, bộ tham số
PI tối ưu, hiệu quả huấn luyện mạng ANN và kết quả đáp ứng hệ thống trên phần cứng
ESP32.


% ==============================================
\section{Kết quả nhận dạng hàm truyền động cơ NF5475}
\label{sec:result_tf}
% ==============================================

Quy trình nhận dạng hàm truyền mô tả trong Mục~\ref{sec:data_sysid} được thực hiện
trên \textbf{sáu mẫu động cơ NF5475 độc lập}. Mục đích thực hiện trên nhiều mẫu nhằm
đánh giá mức độ đồng đều về đặc tính giữa các cá thể động cơ, xác định mẫu đại diện
để thiết kế bộ điều khiển, và phân tích ảnh hưởng của hao mòn cơ học sau thời gian
sử dụng.

% ----------------------------------------------------------
\subsection{Sáu hàm truyền xác định được}
\label{subsec:six_tf}
% ----------------------------------------------------------

Mỗi mẫu động cơ được kích thích bậc thang ở 5 mức PWM, dữ liệu vận tốc ghi qua
encoder 200~PPR. Hàm truyền $G(s)$ được nhận dạng bằng \texttt{tfest} với cấu trúc
\textbf{bậc 4/3} (4 cực, 3 không). Sáu kết quả:

\begin{equation}
    G_1(s) = \frac{353{,}5\,s^3 + 81{,}55\,s^2 + 16{,}27\,s + 3{,}984}
                  {s^4 + 33{,}00\,s^3 + 7{,}444\,s^2 + 3489\,s + 0{,}2802}
    \label{eq:tf1}
\end{equation}

\begin{equation}
    G_2(s) = \frac{361{,}2\,s^3 + 79{,}38\,s^2 + 15{,}84\,s + 4{,}012}
                  {s^4 + 34{,}21\,s^3 + 7{,}512\,s^2 + 3502\,s + 0{,}2756}
    \label{eq:tf2}
\end{equation}

\begin{equation}
    G_3(s) = \frac{348{,}7\,s^3 + 83{,}21\,s^2 + 16{,}58\,s + 3{,}951}
                  {s^4 + 32{,}54\,s^3 + 7{,}398\,s^2 + 3478\,s + 0{,}2831}
    \label{eq:tf3}
\end{equation}

\begin{equation}
    G_4(s) = \frac{355{,}9\,s^3 + 80{,}92\,s^2 + 16{,}43\,s + 4{,}008}
                  {s^4 + 33{,}41\,s^3 + 7{,}487\,s^2 + 3495\,s + 0{,}2814}
    \label{eq:tf4}
\end{equation}

\begin{equation}
    G_5(s) = \frac{359{,}1\,s^3 + 82{,}47\,s^2 + 16{,}11\,s + 3{,}967}
                  {s^4 + 33{,}82\,s^3 + 7{,}456\,s^2 + 3511\,s + 0{,}2789}
    \label{eq:tf5}
\end{equation}

\begin{equation}
    G_6(s) = \frac{352{,}4\,s^3 + 80{,}13\,s^2 + 15{,}97\,s + 3{,}996}
                  {s^4 + 32{,}89\,s^3 + 7{,}431\,s^2 + 3485\,s + 0{,}2823}
    \label{eq:tf6}
\end{equation}

Chỉ số đánh giá của từng hàm truyền và hệ số khuếch đại một chiều (DC gain)
được tổng hợp trong Bảng~\ref{tab:tf_metrics}. Hệ số khuếch đại DC được tính
theo định lý giá trị cuối:
\begin{equation}
    K_{DC} = \lim_{s \to 0} G_i(s)
           = \frac{b_0}{a_0}
    \label{eq:dc_gain}
\end{equation}
trong đó $b_0$ và $a_0$ lần lượt là hệ số tự do của tử và mẫu số.

\begin{table}[H]
    \centering
    \caption{Chỉ số đánh giá sáu hàm truyền động cơ NF5475}
    \label{tab:tf_metrics}
    \renewcommand{\arraystretch}{1.3}
    \begin{tabular}{ccccc}
        \toprule
        \textbf{Mẫu} & \textbf{Độ chính xác (\%)} & \textbf{FPE} & \textbf{MSE} & \textbf{$K_{DC}$} \\
        \midrule
        $G_1(s)$ & 95{,}36 & 2\,572 & 2\,562 & 14{,}22 \\
        $G_2(s)$ & 94{,}87 & 2\,698 & 2\,687 & 14{,}56 \\
        $G_3(s)$ & 93{,}42 & 3\,012 & 2\,997 & 13{,}96 \\
        $G_4(s)$ & 95{,}12 & 2\,614 & 2\,603 & 14{,}24 \\
        $G_5(s)$ & 94{,}23 & 2\,821 & 2\,809 & 14{,}22 \\
        $G_6(s)$ & 94{,}61 & 2\,735 & 2\,724 & 14{,}15 \\
        \midrule
        \textbf{TB $\pm$ std} & $94{,}94 \pm 0{,}65$ & $2\,742 \pm 152$ & $2\,730 \pm 153$ & $14{,}23 \pm 0{,}18$ \\
        \bottomrule
    \end{tabular}
\end{table}

$G_1(s)$ được chọn làm hàm truyền đại diện vì đạt độ chính xác cao nhất
(95,36\%) và FPE/MSE thấp nhất trong nhóm.

% ----------------------------------------------------------
\subsection{Phân tích hàm truyền đại diện $G_1(s)$}
\label{subsec:tf1_analysis}
% ----------------------------------------------------------

\paragraph{Cấu trúc mô hình.}
Hàm truyền $G_1(s)$ có bậc 4 ở mẫu số và bậc 3 ở tử số.
Điều này phản ánh bốn cơ chế động học chồng lên nhau trong hệ động cơ thực:
\begin{enumerate}
    \item \textbf{Động học điện (electrical dynamics):} tương ứng với hằng số
    thời gian điện $\tau_e = L_a/R_a$, đặc trưng cho cực có phần thực lớn (âm, tần số cao).

    \item \textbf{Động học cơ khí (mechanical dynamics):} tương ứng với
    $\tau_m = J/B$, đặc trưng cho cực có phần thực nhỏ hơn.

    \item \textbf{Ma sát phi tuyến:} Hiệu ứng Coulomb và ma sát Stribeck tạo
    ra dao động nhỏ khi thay đổi chiều và ở vận tốc thấp;
    \texttt{tfest} tuyến tính hóa thành cực bổ sung.

    \item \textbf{Động học bộ mã hóa và lọc số:} Encoder quang học 200~PPR
    kết hợp bộ lọc trung bình động (window~$=5$ mẫu, xem Mục~\ref{sec:data_sysid})
    cộng thêm độ trễ tổng hợp vào mô hình.
\end{enumerate}

\paragraph{Hệ số khuếch đại DC.}
\begin{equation}
    K_{DC,1} = \frac{3{,}984}{0{,}2802} \approx 14{,}22 \; \frac{\text{RPM}}{\text{đơn vị PWM}}
\end{equation}
Giá trị này phản ánh độ nhạy vận tốc theo đầu vào PWM ở trạng thái xác lập.
Hệ số $K_{DC}$ của sáu mẫu dao động trong khoảng $13{,}96 \div 14{,}56$,
tức biến động $\approx \pm 2\%$ so với trung bình — xác nhận tính đồng đều của
lô động cơ.

\paragraph{Chất lượng xấp xỉ.}
Độ chính xác 95,36\% của $G_1(s)$ cho biết mô hình tuyến tính bậc 4/3 giải thích
được phần lớn động học thực. Phần sai số còn lại (4,64\%) chủ yếu do:
\begin{itemize}
    \item Ma sát tĩnh và vùng chết (dead-zone) không thể biểu diễn bằng hàm truyền tuyến tính.
    \item Sự thay đổi điện trở phần ứng theo nhiệt độ trong quá trình thực nghiệm.
    \item Nhiễu lượng tử hóa encoder (200~PPR → độ phân giải $\approx 0{,}3\,^\circ$/xung).
\end{itemize}
Giá trị FPE~=~2572 và MSE~=~2562 phản ánh biên độ nhiễu đo lường vận tốc
(đơn vị: $\text{RPM}^2$), phù hợp với tín hiệu vận tốc dao động $\pm 10$~RPM
quanh setpoint trong điều kiện thực nghiệm không phòng sạch.

% ----------------------------------------------------------
\subsection{Nhận xét đặc tính động cơ}
\label{subsec:motor_char}
% ----------------------------------------------------------

\paragraph{Sự phân tán giữa các mẫu.}
Mặc dù sáu mẫu cùng mã NF5475, các hệ số hàm truyền dao động nhỏ nhưng có hệ thống.
Các nguyên nhân chính:

\begin{itemize}
    \item \textbf{Dung sai chế tạo:} Điện trở phần ứng $R_a$, điện cảm $L_a$ và
    hệ số ma sát nhớt $B$ của từng cá thể động cơ không hoàn toàn giống nhau
    ngay từ khi xuất xưởng.

    \item \textbf{Mức độ hao mòn chổi than khác nhau:}
    Động cơ DC có chổi than (brushed DC) sử dụng lâu dài sẽ bị mòn chổi than,
    làm tăng điện trở tiếp xúc $R_c$. Điện trở tổng $R_{total} = R_a + R_c$
    tăng dần, kéo theo hằng số thời gian điện $\tau_e = L_a/R_{total}$ giảm
    và DC gain $K_{DC} \propto 1/R_{total}$ giảm nhẹ. Đây là lý do $G_3(s)$
    có $K_{DC}$ thấp nhất (13,96) — khả năng mẫu này đã qua sử dụng nhiều hơn.

    \item \textbf{Bụi carbon từ mòn chổi:}
    Bụi carbon bám trên cổ góp (commutator) gây dao động tiếp xúc ngẫu nhiên,
    làm tăng FPE/MSE. $G_3(s)$ có FPE cao nhất (3012), nhất quán với giả thuyết
    mẫu này có chổi than mòn nhiều hơn.

    \item \textbf{Mòn ổ bi và tăng ma sát Coulomb:}
    Ổ bi mòn tăng lực cản cơ học $B$, rút ngắn hằng số thời gian cơ $\tau_m = J/B$.
    Điều này thể hiện qua sự thay đổi nhỏ ở hệ số mẫu số bậc cao
    ($33{,}00 \div 34{,}21$ cho hệ số $s^3$).

    \item \textbf{Suy giảm từ trường nam châm vĩnh cửu:}
    Với động cơ PMDC, từ thông $\Phi$ giảm nhẹ theo tuổi thọ (nhiệt độ cao,
    rung động), làm giảm hằng số mômen $K_m$ và hằng số sức phản điện động
    $K_e$, ảnh hưởng đến hệ số khuếch đại vòng hở.
\end{itemize}

\paragraph{Khuyến nghị cho thiết kế bộ điều khiển.}
Do sự phân tán giữa các mẫu ($K_{DC}$ biến động $\approx 4\%$, FPE biến động
$\approx 17\%$), bộ điều khiển được thiết kế trên $G_1(s)$ cần có \textbf{biên
pha và biên độ đủ lớn} để hoạt động ổn định trên toàn bộ lô động cơ.
Đây là một trong các lý do sử dụng giải thuật di truyền (GA) thay vì thiết kế
phân tích trực tiếp: GA cho phép điều chỉnh ngưỡng phạt tối ưu theo đặc tính
thực đo của từng cá thể, thay vì giả định mô hình lý tưởng.


% ==============================================
\section{Kết quả tối ưu tham số PI bằng GA}
\label{sec:result_ga}
% ==============================================

% ----------------------------------------------------------
\subsection{So sánh phương pháp tối ưu}
\label{subsec:compare_methods}
% ----------------------------------------------------------

Ba phương pháp tối ưu bộ điều khiển PI vận tốc được so sánh trên cùng mô hình
$G_1(s)$ rời rạc hóa (Tustin, $T_s = 0{,}01\,$s):

\begin{itemize}
    \item \textbf{Ngẫu nhiên trong miền (Random-in-bounds):} Lấy mẫu ngẫu nhiên
    đều trong không gian tham số $K_p \in [0{,}01,\,1]$, $K_i \in [0{,}1,\,10]$,
    chọn nghiệm tốt nhất sau 1\,000 phép thử.

    \item \textbf{MATLAB PID Tuner:} Sử dụng công cụ tự động của MATLAB dựa trên
    phương pháp phân tích vòng hở (loopshaping), chỉnh theo tiêu chí thời gian
    đáp ứng mong muốn.

    \item \textbf{Giải thuật di truyền (GA):} Cấu hình theo Mục~\ref{sec:ga_optim}:
    80 cá thể, 80 thế hệ, 20 lần chạy độc lập, hàm mục tiêu có hàm phạt
    (overshoot $> 5\%$ và thời gian xác lập $> 0{,}5\,$s).
\end{itemize}

\begin{table}[H]
    \centering
    \caption{So sánh ba phương pháp tối ưu tại setpoint 250~RPM}
    \label{tab:method_compare}
    \renewcommand{\arraystretch}{1.3}
    \begin{tabular}{lcccc}
        \toprule
        \textbf{Phương pháp} & \textbf{Overshoot (\%)} & \textbf{$t_{s}$ (s)} & \textbf{$t_{r}$ (s)} & \textbf{SSE (RPM)} \\
        \midrule
        Random-in-bounds & 18{,}4 & 1{,}32 & 0{,}28 & 12{,}7 \\
        MATLAB PID Tuner & 8{,}2  & 0{,}74 & 0{,}21 & 3{,}4  \\
        GA (đề xuất)     & 3{,}2  & 0{,}38 & 0{,}19 & 0{,}8  \\
        \bottomrule
    \end{tabular}
\end{table}

Trong đó $t_s$ là thời gian xác lập (sai số $< 2\%$), $t_r$ là thời gian lên,
SSE là sai số xác lập trung bình. GA vượt trội cả ba chỉ tiêu,
đặc biệt giảm overshoot từ 18,4\% (ngẫu nhiên) xuống còn 3,2\% và sai số
xác lập giảm 15,9 lần so với phương pháp ngẫu nhiên.

% ----------------------------------------------------------
\subsection{Bộ gain tối ưu tại 9 mức setpoint}
\label{subsec:ga_nine_setpoints}
% ----------------------------------------------------------

Quy trình GA được lặp lại cho \textbf{9 mức setpoint vận tốc} từ 250~RPM đến
2250~RPM (bước nhảy 250~RPM). Với mỗi setpoint, chỉ hai tham số vòng trong
$(K_{pv}^*, K_{iv}^*)$ được trích xuất để dùng cho ANN (xem Mục~\ref{sec:ann_train});
ba tham số vòng ngoài $(K_{pp}^*, K_{ip}^*, K_{dp}^*)$ giữ nguyên.

\begin{table}[H]
    \centering
    \caption{Kết quả tối ưu GA: bộ tham số PI vận tốc tại 9 setpoint}
    \label{tab:ga_nine}
    \renewcommand{\arraystretch}{1.3}
    \begin{tabular}{ccccc}
        \toprule
        \textbf{Setpoint (RPM)} & $K_{pv}^*$ & $K_{iv}^*$ & \textbf{Overshoot (\%)} & \textbf{$t_s$ (s)} \\
        \midrule
        250  & 0{,}1854 & 5{,}288 & 3{,}2 & 0{,}38 \\
        500  & 0{,}1721 & 5{,}104 & 3{,}8 & 0{,}42 \\
        750  & 0{,}1698 & 4{,}932 & 4{,}1 & 0{,}45 \\
        1000 & 0{,}1612 & 4{,}847 & 4{,}3 & 0{,}48 \\
        1250 & 0{,}1583 & 4{,}721 & 4{,}5 & 0{,}51 \\
        1500 & 0{,}1547 & 4{,}658 & 4{,}7 & 0{,}53 \\
        1750 & 0{,}1512 & 4{,}512 & 4{,}8 & 0{,}56 \\
        2000 & 0{,}1487 & 4{,}389 & 4{,}9 & 0{,}59 \\
        2250 & 0{,}1463 & 4{,}274 & 5{,}0 & 0{,}62 \\
        \bottomrule
    \end{tabular}
\end{table}

Mỗi bộ $\{K_{pv}^*, K_{iv}^*\}$ được kiểm chứng trên dữ liệu thực nghiệm
10~giây (1001~mẫu) và đảm bảo: overshoot $\leq 5\%$, $t_s \leq 0{,}65\,$s.

% ----------------------------------------------------------
\subsection{Phân tích xu hướng tham số}
\label{subsec:ga_trend}
% ----------------------------------------------------------

Từ Bảng~\ref{tab:ga_nine}, có thể quan sát hai xu hướng rõ ràng:

\begin{enumerate}
    \item \textbf{$K_{pv}^*$ giảm dần khi setpoint tăng}
    (từ 0,1854 ở 250~RPM xuống 0,1463 ở 2250~RPM, giảm $\approx 21\%$).
    Ở vận tốc cao, sức phản điện động (back-EMF) $e = K_e \omega$ tăng lên,
    làm tăng hệ số khuếch đại vòng hở hiệu dụng. Bộ điều khiển cần giảm
    $K_{pv}$ để duy trì biên pha ổn định và tránh vọt lố.

    \item \textbf{$K_{iv}^*$ giảm dần khi setpoint tăng}
    (từ 5,288 ở 250~RPM xuống 4,274 ở 2250~RPM, giảm $\approx 19\%$).
    Ở vận tốc cao, nhiễu đo lường encoder tỉ lệ với vận tốc nên tích phân
    $K_{iv}$ lớn có thể khuếch đại nhiễu và gây dao động. GA tự động tìm
    ra giá trị $K_{iv}$ thấp hơn để cân bằng giữa triệt tiêu sai số xác
    lập và độ ổn định.
\end{enumerate}

Xu hướng đơn điệu và nhất quán của $\{K_{pv}^*, K_{iv}^*\}$ theo setpoint
tạo điều kiện thuận lợi để mạng ANN khái quát hóa (generalize) tốt, vì
hàm ánh xạ setpoint $\to$ tham số là gần tuyến tính và đơn điệu
(xem Mục~\ref{sec:ann_train}).

\paragraph{Nhận xét về sự hội tụ của GA.}
Với cấu hình 80 cá thể, 80 thế hệ và 20 lần chạy độc lập:
\begin{itemize}
    \item Độ lệch chuẩn nghiệm tốt nhất giữa 20 lần chạy $< 2\%$,
    cho thấy GA đã hội tụ ổn định và không bị kẹt tại cực tiểu địa phương.
    \item Thời gian chạy trung bình với \texttt{UseParallel = true}
    (4 luồng): $\approx 45$ giây/setpoint trên CPU Intel Core i7.
    \item Tổng thời gian tối ưu cho 9 setpoint: $\approx 6{,}8$ phút.
\end{itemize}


% ==============================================
\section{Kết quả huấn luyện mạng ANN}
% ==============================================

\subsection{Tóm tắt kiến trúc và cấu hình}

Kiến trúc và quy trình huấn luyện được trình bày chi tiết trong
Chương~\ref{chap:phuongphap}. Tóm tắt:
\begin{itemize}
    \item Kiến trúc: $2 \to 8 \to 16 \to 2$ (Dense + L2 + ReLU + Dropout).
    \item Optimizer: Adam, lr $= 2\times10^{-3}$, decay $= 3\times10^{-4}$.
    \item Số epochs: 1\,000.
\end{itemize}

\subsection{Kết quả huấn luyện}

Mạng ANN hội tụ sau khoảng 400 epochs. Đường cong loss huấn luyện và validation
bám sát nhau, chứng tỏ mạng không bị overfitting nhờ L2 regularization và Dropout
(xem Hình~\ref{fig:loss_curve}).

\begin{figure}[H]
    \centering
    % \includegraphics[width=0.75\textwidth]{figure/training_loss.png}
    \caption{Đường cong Training Loss và Validation Loss theo epoch}
    \label{fig:loss_curve}
\end{figure}

\subsection{Đánh giá mô hình}

Mô hình được đánh giá trên tập test (2 bộ dữ liệu chưa dùng để huấn luyện).
Kết quả dự đoán $K_p$, $K_i$ so với giá trị tối ưu từ GA cho thấy sai số nhỏ
và ổn định trên toàn dải setpoint 250 -- 2250 RPM.

\textit{Hạn chế hiện tại:} Chưa có số liệu định lượng đầy đủ về độ chính xác dự đoán
trên phần cứng thực; đây là hướng cần bổ sung trong giai đoạn tiếp theo.


% ==============================================
\section{Kết quả huấn luyện chính sách PPO trong Isaac Lab}
\label{sec:result_rl}
% ==============================================

% ----------------------------------------------------------
\subsection{Tổng quan quá trình huấn luyện}
\label{subsec:rl_overview}
% ----------------------------------------------------------

Quá trình huấn luyện được thực hiện trên GPU NVIDIA GeForce RTX~3070 (8~GB VRAM)
với 4\,096 môi trường song song, tổng cộng 5\,000 lần lặp.
Hình~\ref{fig:rl_overview} trình bày hai chỉ số tổng quan xuyên suốt quá trình:
phần thưởng trung bình episode (\textit{mean\_reward}) và độ dài episode trung
bình (\textit{mean\_episode\_length}).

\begin{figure}[H]
    \centering
    \includegraphics[width=\textwidth]{figure/results/01_train_overview.png}
    \caption{Tổng quan huấn luyện: phần thưởng episode và độ dài episode trung bình
             theo số lần lặp}
    \label{fig:rl_overview}
\end{figure}

Có thể phân tách quá trình huấn luyện thành ba giai đoạn rõ rệt:

\begin{enumerate}
    \item \textbf{Giai đoạn khởi động (iter 0-1\,400):}
    Phần thưởng tăng dần từ $\approx 0$ lên $\approx 170$; độ dài episode tăng
    từ 0 lên $\approx 1\,700$ bước ($\approx 57\%$ chiều dài tối đa 3\,000 bước).
    Robot học giữ thăng bằng cơ bản nhưng vẫn thường xuyên ngã.

    \item \textbf{Giai đoạn mất ổn định (iter 1\,400--1\,900):}
    Xuất hiện suy giảm đột ngột: độ dài episode giảm từ $\approx 1\,700$ xuống
    $\approx 800$ bước, phần thưởng giảm về $\approx 100$. Giai đoạn này tương
    ứng với sự tăng mạnh tỉ lệ tiếp xúc bất hợp lệ chân trái (xem
    Hình~\ref{fig:rl_terminations}), cho thấy chính sách đang trải qua tái
    tổ chức hành vi trước khi hội tụ.

    \item \textbf{Giai đoạn hội tụ (iter 1\,900--4\,800):}
    Phần thưởng tăng nhanh từ $\approx 100$ lên $\approx 650$; độ dài episode
    vọt lên $\approx 2\,900$ bước ($\approx 97\%$ chiều dài tối đa) và duy trì
    ổn định. Robot thành thạo nhiệm vụ Phase~0: đứng thẳng, duy trì chiều cao
    và cân bằng thân trong toàn bộ episode 60~giây.
\end{enumerate}

\textbf{Model tốt nhất được chọn tại checkpoint iter~4\,800} (lưu theo chu kỳ
100 lần lặp), nơi phần thưởng đạt đỉnh $\approx 650$.

% ----------------------------------------------------------
\subsection{Phân tích phần thưởng thành phần}
\label{subsec:rl_rewards}
% ----------------------------------------------------------

\paragraph{Phần thưởng nhiệm vụ (Task Rewards).}

\begin{figure}[H]
    \centering
    \includegraphics[width=\textwidth]{figure/results/02_task_rewards.png}
    \caption{Diễn biến các phần thưởng nhiệm vụ theo số lần lặp}
    \label{fig:rl_task_rewards}
\end{figure}

Tại iter~4\,800 (model tốt nhất), giá trị các phần thưởng nhiệm vụ:

\begin{table}[H]
    \centering
    \caption{Giá trị phần thưởng nhiệm vụ tại checkpoint iter~4\,800}
    \label{tab:task_reward_final}
    \renewcommand{\arraystretch}{1.3}
    \begin{tabular}{lccc}
        \toprule
        \textbf{Thành phần} & \textbf{Trọng số} & \textbf{Giá trị đạt được} & \textbf{Hiệu suất} \\
        \midrule
        \texttt{track\_lin\_vel\_xy\_exp} & $+5{,}0$ & $\approx 4{,}8$ & $96\%$ \\
        \texttt{track\_ang\_vel\_z\_exp}  & $+4{,}0$ & $\approx 3{,}1$ & $78\%$ \\
        \texttt{track\_base\_height\_exp} & $+8{,}0$ & $\approx 4{,}0$ & $50\%$ \\
        \texttt{upright}                  & $-5{,}0$ & $\approx -0{,}1$ & $98\%$\textsuperscript{*} \\
        \bottomrule
        \multicolumn{4}{l}{\small \textsuperscript{*}Hiệu suất upright = tỉ lệ gần 0 so với mức phạt tệ nhất $-5{,}0$}
    \end{tabular}
\end{table}

\textbf{Nhận xét:}
Lệnh vận tốc trong toàn bộ lần chạy luôn là $v_x^{cmd} = 0$, $\omega_z^{cmd} = 0$
% (Phase~0, xem Mục~\ref{subsec:rl_curriculum}). Do đó \texttt{track\_lin\_vel\_xy\_exp}
đạt 96\% đồng nghĩa với việc robot đứng gần yên hoàn toàn ($v_x \approx 0$~m/s).
Phần thưởng \texttt{upright} gần bằng 0 cho thấy thân robot duy trì tư thế thẳng
đứng rất tốt. Phần thưởng chiều cao đạt 50\% phản ánh robot kiểm soát được chiều
cao thân nhưng còn dao động quanh setpoint. Phần thưởng vận tốc quay đạt 78\%
tuy nhiên có yaw drift nhỏ do không có cơ chế chủ động chỉnh hướng khi đứng yên.

\paragraph{Phần thưởng phạt (Penalty Rewards).}

\begin{figure}[H]
    \centering
    \includegraphics[width=\textwidth]{figure/results/03_penalty_rewards.png}
    \caption{Diễn biến các phần thưởng phạt theo số lần lặp}
    \label{fig:rl_penalty_rewards}
\end{figure}

Tại iter~4\,800:
\begin{itemize}
    \item \texttt{action\_rate} $\approx -0{,}45$: Robot tạo ra thay đổi hành
    động đáng kể giữa các bước, phản ánh hoạt động tích cực của khớp hông để
    duy trì thăng bằng động.

    \item \texttt{joint\_deviation\_pad} $\approx -0{,}35$: Khớp hông duy trì
    lệch khỏi vị trí mặc định ở mức vừa phải — đây là hành vi mong muốn vì
    robot cần điều chỉnh chủ động tư thế.

    \item \texttt{joint\_acc} $\approx -0{,}02$, \texttt{joint\_torques} $\approx -0{,}07$:
    Rất nhỏ, cho thấy chuyển động khớp mượt mà và mô-men xoắn thấp.

    \item \texttt{termination\_penalty} $\approx 0$: Hầu như không có termination
    sớm — nhất quán với việc robot hoàn thành toàn bộ episode.
\end{itemize}

% ----------------------------------------------------------
\subsection{Chỉ số bám lệnh và điều kiện kết thúc}
\label{subsec:rl_tracking}
% ----------------------------------------------------------

\begin{figure}[H]
    \centering
    \includegraphics[width=\textwidth]{figure/results/04_tracking_metrics.png}
    \caption{Chỉ số sai số bám lệnh theo số lần lặp}
    \label{fig:rl_tracking}
\end{figure}

\texttt{error\_vel\_xy} trong giai đoạn đầu (iter 0--1\,400) tăng lên $\approx 11$:
robot chưa học đứng yên khi lệnh vận tốc bằng 0, bánh xe vẫn quay không kiểm soát.
Sau iter~1\,900, \texttt{error\_vel\_xy} giảm nhanh và ổn định ở $\approx 0{,}7$
tức nghĩa là sai số vận tốc dọc rất nhỏ, robot đứng gần yên hoàn toàn.
\texttt{error\_vel\_yaw} ổn định ở $\approx 3{,}0$ từ iter~2\,000 còn drift
nhỏ theo trục yaw nhưng không gây mất ổn định.
\texttt{position\_error} và \texttt{orientation\_error} đều xấp xỉ 0 trong toàn
bộ quá trình, xác nhận robot duy trì tư thế thẳng đứng tốt.

\begin{figure}[H]
    \centering
    \includegraphics[width=\textwidth]{figure/results/05_terminations.png}
    \caption{Tỉ lệ các điều kiện kết thúc episode theo số lần lặp}
    \label{fig:rl_terminations}
\end{figure}

Điều kiện kết thúc cho thấy quá trình học rõ ràng:
\begin{itemize}
    \item \texttt{illegal\_contact\_left}: khởi đầu $\approx 95\%$ (robot thường
    xuyên đổ về phía trái), giảm dần và đạt $\approx 1\%$ sau iter~2\,500.
    Robot đã học được phối hợp tay chân để tránh tiếp xúc bất hợp lệ.

    \item \texttt{illegal\_contact\_right}: bắt đầu $\approx 30\%$, giảm về
    $\approx 2\%$ sau iter~2\,000. Thấp hơn bên trái do bất đối xứng ngẫu nhiên
    trong khởi tạo.

    \item \texttt{time\_out}: tăng từ 0 lên $\approx 95\%$ sau iter~2\,500 và
    duy trì ổn định tận \textbf{95\% episode hoàn thành đủ 60 giây mà không ngã}.
    Đây là chỉ số quan trọng nhất khẳng định robot thành thạo nhiệm vụ cân bằng.

    \item Các điều kiện còn lại (\texttt{bad\_orientation}, \texttt{illegal\_contact\_base},
    \texttt{joint\_vel\_limit}) duy trì gần 0 trong toàn bộ quá trình.
\end{itemize}

% ----------------------------------------------------------
\subsection{Phân tích PPO Loss}
\label{subsec:rl_loss}
% ----------------------------------------------------------

\begin{figure}[H]
    \centering
    \includegraphics[width=\textwidth]{figure/results/06_loss.png}
    \caption{Các thành phần loss của thuật toán PPO theo số lần lặp}
    \label{fig:rl_loss}
\end{figure}

\begin{itemize}
    \item \textbf{Surrogate loss} (xanh lam): duy trì $\approx 0$ suốt quá trình
    — cơ chế clipping PPO ($\varepsilon = 0{,}2$) hoạt động đúng, tránh cập nhật
    chính sách quá lớn.

    \item \textbf{Value function loss} (cam): bắt đầu $\approx 2{,}4$, giảm về
    $\approx 0$ sau iter~4\,000 — Critic học dự đoán chính xác giá trị hàm
    return. Có nhồng nhỏ ở iter $\approx 1\,500$ (tương ứng giai đoạn mất ổn định)
    trước khi tiếp tục giảm.

    \item \textbf{Entropy} (xanh lá): giảm từ $\approx 2{,}3$ xuống $\approx -12$
    đơn điệu trong 5\,000 iter — chính sách chuyên hóa dần và trở nên
    deterministic hơn. Entropy âm (log-entropy) là bình thường với phân phối
    Gaussian hẹp khi chính sách hội tụ.
\end{itemize}

% ----------------------------------------------------------
\subsection{Lựa chọn model tốt nhất}
\label{subsec:rl_best_model}
% ----------------------------------------------------------

Dựa trên đường cong phần thưởng (Hình~\ref{fig:rl_overview}), phần thưởng
episode đạt đỉnh $\approx 650$ tại iter $\approx 4\,800$ trước khi có sự sụt
giảm nhỏ ở cuối. \textbf{Checkpoint tại iteration~4\,800} được chọn làm model
tốt nhất cho giai đoạn triển khai tiếp theo.

\begin{table}[H]
    \centering
    \caption{Tổng kết hiệu năng model tốt nhất (iter~4\,800)}
    \label{tab:best_model}
    \renewcommand{\arraystretch}{1.3}
    \begin{tabular}{lc}
        \toprule
        \textbf{Chỉ số} & \textbf{Giá trị} \\
        \midrule
        Phần thưởng episode trung bình                  & $\approx 650$ \\
        Độ dài episode trung bình                       & $\approx 2\,900$ / 3\,000 bước ($97\%$) \\
        Tỉ lệ hoàn thành đủ 60~s (\texttt{time\_out})  & $\approx 95\%$ \\
        Sai số vận tốc dọc (\texttt{error\_vel\_xy})    & $\approx 0{,}7$ \\
        Sai số vận tốc quay (\texttt{error\_vel\_yaw})  & $\approx 3{,}0$ \\
        Phần thưởng thẳng đứng (\texttt{upright})      & $\approx -0{,}1$ / $-5{,}0$ (tệ nhất) \\
        Phần thưởng bám vận tốc dọc                    & $\approx 4{,}8$ / $5{,}0$ ($96\%$) \\
        Phần thưởng bám chiều cao                      & $\approx 4{,}0$ / $8{,}0$ ($50\%$) \\
        Value function loss                            & $\approx 0$ (hội tụ) \\
        \bottomrule
    \end{tabular}
\end{table}

Trọng số Actor tại iter~4\,800 được xuất sang định dạng \textbf{ONNX}
(\textit{Open Neural Network Exchange}) để chuẩn bị triển khai lên phần cứng
thực tế trong giai đoạn tiếp theo.
